\section{Контекст до миллиона токенов}
\label{sec:million}

Чем глубже контекст, тем больше агент держит в голове: весь репозиторий, длинную историю работы, большой набор
документов. Налог на поддержание кэша растёт вместе с глубиной, потому что каждый вызов перечитывает весь контекст. У
поставщика GLM-5.3-Flash сессия, которая дорастает до 200 тысяч токенов, стоит \$0.44, а сессия до миллиона \$9.15:
контекст глубже в пять раз, сессия дороже в 21 раз.

\textbf{Клиент для миллиона токенов.} GLM-5.3-Flash устроена как MoE: 328 миллиардов параметров, из них 17.5 активных
\cite{glm53flash}. Её кэш внимания занимает около 6~КБ на токен, в шесть раз меньше, чем у Qwen3.8-27B, а рекуррентное состояние около 153~МБ в памяти и 76~МБ при передаче в 16 битах. Веса в 4 битах занимают 181~ГБ и помещаются в Mac Studio с M5 Ultra и 256~ГБ памяти. Рядом с весами
помещается 10 сессий по миллиону токенов или 46 сессий по 200 тысяч. Генерация, по нашей оценке, идёт со скоростью
около 19 токенов в секунду на поток при четырёх потоках. Узкое место то же, что и у Qwen3.8-27B: прочитать миллион токенов сама Studio может по нашей оценке примерно за 57 минут, а одна B300 считает их prefill за
43 секунды, и вместе с передачей по гигабитному каналу выходит около 97 секунд. Qwen3.8-27B
на клиенте с 128~ГБ по объёму памяти держит две сессии по миллиону токенов, но это за пределами её штатного
контекста (таблица~\ref{tab:fit}).

\textbf{Потеря кэша.} Клиент, который потерял и память, и сохранённое на диске состояние сессии в миллион токенов,
почти теряет саму сессию. M3 Max читал бы такую сессию Qwen3.8-27B заново 26--28 часов: время на токен по замеренным
отрезкам prefill (таблица~\ref{tab:m3max}) мы продолжаем линейно за 231 тысячу токенов с наклоном отрезка от 181 до
231 тысячи (25.9 часа) или от 131 до 231 тысячи (28.2 часа). Это оценка за пределами замеров и штатного контекста
модели. В гибриде сервер хранит копию сессии и присылает клиенту её память целиком: на миллионе токенов это около
35~ГБ, по гигабитному каналу около пяти минут, трафик по \$5 за терабайт стоит около 18 центов. Хранение копии на NVMe по
\$0.10 за гигабайт стоит около \$3.5 в месяц. По политике протокола копия в рабочие часы лежит в оперативной памяти
по \$0.50 за гигабайт, и хранение стоит около \$6.7. Пять минут здесь нижняя граница: по политике протокола
сервер дописывает в копию токены вызовов, прочитанных клиентом, только при следующем оффлоаде, и перед
восстановлением ему может понадобиться prefill этого хвоста. Если диск клиента уцелел, сессию быстрее поднять из
локального снимка на его SSD.

\textbf{Стоимость сессии.} Рисунок~\ref{fig:million} сравнивает стоимость сессии до миллиона токенов у API, у своего
серверного кластера и у гибрида с Studio. Гибрид включает покупку Studio целиком: \$10\,799 с амортизацией за четыре
года при остаточной стоимости 35\% и электричеством дают \$149 в месяц. При двух таких сессиях в день на человека это \$3.54 на сессию только за Studio. Для команд от
500 человек гибрид обходится в \$4.2--4.6 за сессию против \$9.15 у API, а свой кластер стоит \$3.21 и дешевле
гибрида.
Команде из 10 человек и меньше выделенный сервер не окупается, и дешевле всего остаётся API. На 50 человек свой кластер
(\$8.02) немного дешевле API (\$9.15), а гибрид с отдельной Studio на каждого дороже обоих (\$11.85). Малым командам нужен общий
поставщик prefill.

Studio при двух сессиях в день загружена меньше чем на треть: в гибридном режиме она успевает провести почти семь
сессий по миллиону токенов за рабочий день. При полной загрузке доля Studio падает до \$1.04 на сессию, и сессия
гибрида для команды от 500 человек стоит \$1.5--1.7, в 5.4--6.0 раза дешевле API и в 2.0--2.1 раза дешевле своего
кластера. Четыре сессии в день на человека дают \$2.3--2.4, в 1.4 раза дешевле своего кластера; если Studio делят двое и каждый проводит по две сессии в день,
\$2.4--2.9. При осторожной политике хранения сессии дороже не больше чем на 20\%. Загрузка клиента и есть главный
рычаг гибрида на глубоком контексте.

\begin{figure}[t]
\centering
\begin{tikzpicture}
\begin{axis}[
  width=0.62\textwidth, height=55mm, xmode=log, ymode=log,
  xmin=4, xmax=6500, ymin=2, ymax=120, xtick={5,10,20,50,500,5000}, xticklabels={5,10,20,50,500,5000},
  ytick={2,5,10,20,50,100}, yticklabels={2,5,10,20,50,100},
  xlabel={человек в команде}, ylabel={\$ за сессию},
  tick label style={font=\footnotesize}, label style={font=\footnotesize},
  legend pos=north east, legend style={font=\footnotesize,draw=none,fill=none}]
\addplot[black,dashed] table[x=team,y=api,col sep=comma] {data/glm-1m-cost.csv};
\addlegendentry{API}
\addplot[black,mark=*,mark size=1.4pt] table[x=team,y=own,col sep=comma] {data/glm-1m-cost.csv};
\addlegendentry{свой кластер B300}
\addplot[black,very thick] table[x=team,y=hybrid,col sep=comma] {data/glm-1m-cost.csv};
\addlegendentry{гибрид, Studio куплена}
\end{axis}
\end{tikzpicture}
\caption{Стоимость сессии GLM-5.3-Flash, которая дорастает до миллиона токенов, в зависимости от размера команды. Две
такие сессии в день на человека; гибрид включает покупку Mac Studio M5 Ultra 256~ГБ и серверный prefill на своих узлах
B300.}
\label{fig:million}
\end{figure}

\textbf{Вывод.} Контекст до миллиона токенов выгоден гибриду, когда клиент загружен: при четырёх сессиях в день на
Studio гибрид в 1.4 раза дешевле своего кластера и почти в четыре раза дешевле API, при полной загрузке в 2.0--2.1 и
в 5.4--6.0 раза. При двух сессиях в день гибрид дороже своего кластера на 29--44\%, при одной в 2.5--2.9 раза;
малым командам без общего поставщика prefill выгоднее API.
